\scp{\tsc{Southwest Maluku}}{04-06}
The \tsc{Southwest Maluku} node contains the languages of the
geographic region of southwest Maluku (except Wetar Island,
see \srf{04-05}), as well as Makuva on the eastern tip of Timor.
This node is defined on the basis of the changes affecting consonant
final words -- more specifically VCVC-final words,
which have resulted in reflexes in the modern languages showing
VC > CV final metathesis in some environments.
%when compared with their PMP etyma.
Examples are given in \trf{04-25}.

\begin{table}[p]\centering\tcp{\tsc{Southwest Maluku} *VC{\#} > CV{\#}}{04-25}
\begin{adjustbox}{width=\textwidth}
	\begin{threeparttable}\st{2.5pt}
	\begin{tabular}{r|ll@{ }llll}\lsptoprule
*gloss	&	sky	&	wind	&	moon	&	hot	&	inside	&	path\\
PMP	&	*laŋ\tbr{it}	&	*haŋ\tbr{in}	&	*bul\tbr{an}	&	*ma-pan\tbr{as}	&	*dal\tbr{əm}	&	*zal\tbr{an}\\	\midrule
Makuva	&		&	{\hr}a\tbr{ne}\su{†}	&	{\hr}hu\tbr{na}\su{†}	&	{\hr}man\tbr{ta}	&		&	{\hr}ʤa\tbr{na}\su{†}\\
\multirow{2}{*}{$\left.\hspace{-2mm}\begin{array}{r}{\textrm{E. Damer}}\\{\textrm{(Ilih)}}\\\end{array}\hspace{-1.3mm}\right\}$}	&	\multirow{2}{*}{{\hr}lian\tbr{ti}}	&	\multirow{2}{*}{{\hr}an\tbr{ni}, anin}	&	\multirow{2}{*}{{\hr}bul\tbr{lo}}	&	\multirow{2}{*}{\cg{(netal)}}	&	\multirow{2}{*}{{\hr}ral\tbr{ma}-n}	&	{\hr}lal\tbr{lo},\\
	&	&	&	&	&	&	{\hr}lalan\\
\multirow{2}{*}{$\left.\hspace{-2mm}\begin{array}{r}{\textrm{E. Damer}}\\{\textrm{(Wulur)}}\\\end{array}\hspace{-1.3mm}\right\}$}	&	\multirow{2}{*}{{\hr}lian\tbr{ti}}	&	\multirow{2}{*}{{\hr}anin}	&	\multirow{2}{*}{{\hr}wul\tbr{no}}	&	\multirow{2}{*}{{\hr}panas\su{‡}}	&	\multirow{2}{*}{{\hr}ral\tbr{ma}-no}	&	{\hr}lalno\\
	&	&	&	&	&	&	\\
Roma	&	{\hr}lyan\tbr{ti}	&	{\hr}an\tbr{ni}	&	{\hr}wul\tbr{la}	&	{\hr}mahan	&	{\hr}ryaram	&	{\hr}tyal\tbr{la}\\
Kisar	&	\cg{(aʔam)}	&	{\hr}an\tbr{ne}	&	{\hr}wol\tbr{lo}	&	{\hr}man\tbr{ha}	&	{\hr}raram	&	{\hr}kal\tbr{la}\\
Luang	&	{\hr}lyan\tbr{ti}	&	{\hr}an\tbr{ni}	&	{\hr}wol\tbr{lə}	&	{\hr}pa\tbr{h}n\tbr{ə}	&	{\hr}ryal\tbr{mə}	&	{\hr}tal\tbr{lə}\\
\multirow{2}{*}{$\left.\hspace{-2mm}\begin{array}{r}{\textrm{Wetan}}\\{\textrm{(Teun Is.)}}\\\end{array}\hspace{-1.3mm}\right\}$}	&	\multirow{2}{*}{{\hr}lian\tbr{ti}}	&	\multirow{2}{*}{{\hr}an\tbr{ni}}	&	\multirow{2}{*}{{\hr}wol\tbr{na}}	&	\multirow{2}{*}{{\hr}pana}	&	\multirow{2}{*}{{\hr}ral\tbr{ma}}	&	\multirow{2}{*}{{\hr}tal\tbr{na}}\\
		&	&	&	&	&	&	\\
\multirow{2}{*}{$\left.\hspace{-2mm}\begin{array}{r}{\textrm{Wetan}}\\{\textrm{(Wetan Is.)}}\\\end{array}\hspace{-1.3mm}\right\}$}	&	\multirow{2}{*}{{\hr}lian\tbr{ti}}	&	\multirow{2}{*}{{\hr}an\tbr{ni}}	&	\multirow{2}{*}{{\hr}wol\tbr{la}}	&	\multirow{2}{*}{{\hr}paana}	&		&	\multirow{2}{*}{{\hr}tal\tbr{la}}\\
		&	&	&	&	&	&	\\
Teun	&	{\hr}lak\tbr{ti}	&	{\hr}wan\tbr{ni}, wanan	&	{\hr}ful\tbr{la}	&	{\hr}pan\tbr{sa}	&	{\hr}ral\tbr{ma}	&	{\hr}tal\tbr{la}\\
Nila	&	{\hr}lak\tbr{ti}	&	{\hr}wan\tbr{ni}, wanin	&	{\hr}ful\tbr{na}	&	{\hr}pan\tbr{sa}	&	{\hr}ral\tbr{na}	&	{\hr}sal\tbr{na}\\
Serua	&	{\hr}lak\tbr{ti}	&	\cg{(motra)}	&	{\hr}wul\tbr{na}	&	\cg{(ni-kreʔu)}	&	{\hr}ral\tbr{na}	&	{\hr}sal\tbr{na}\\
W. Damer	&	\cg{(muhl-o)}	&	{\hr}onyn-o	&	{\hr}ull-o-ni	&	{\hr}hepond-o	&	{\hr}rolm-o	&	{\hr}holl-o\\
Dawera	&	\cg{(dart-ol)}	&	{\hr}aili-ol	&	{\hr}wull-ol	&	\cg{(wetul)}	&	{\hr}ralm	&	{\hr}all-ol\\
Dai	&	\cg{(dart-on)}	&	\cg{(suʔ-on)}	&	{\hr}wunn-on	&	\cg{(wetun)}	&	{\hr}ranm	&	{\hr}hann-on\\
N. Babar	&		&	{\hr}an\tbr{ni}-ai	&	{\hr}wo\tbr{na}	&	\cg{(n-wetunni)}	&	{\hr}rann-e	&	{\hr}haʔan-ai\\
C. Marsela	&	{\hr}lalky-on	&	{\hr}enn-on	&	{\hr}woll	&	{\hr}pant-i\su{Δ}	&	{\hr}ram-n	&	{\hr}all\\
W. Marsela	&	{\hr}lal\tbr{ky}	&	{\hr}unn	&	{\hr}woll	&	{\hr}pant-i	&		&	\\
E.Marsela	&	\cg{(rowek)}	&	{\hr}un	&	{\hr}wol	&	{\hr}pal-pant	&	{\hr}ramn-e	&	{\hr}al\\
Serili	&	{\hr}lalk-ol	&	{\hr}un-en	&	{\hr}bol	&	{\hr}pa-pant-ɛ	&	{\hr}ramn-ɛ	&	{\hr}hal-el\\
SE. Babar	&	{\hr}lalk	&	{\hr}un	&	{\hr}wol	&	{\hr}pal-pan	&	{\hr}ralm	&	{\hr}al\\
Emplawas	&	{\hr}lal\tbr{ky}-o	&	{\hr}u\tbr{sy}-o\su{Δ}	&	{\hr}wos\su{‡}	&	{\hr}pant, pat	&	{\hr}raml-o	&	{\hr}as-o\su{Δ}\\
Tela	&	{\hr}nan\tbr{ky}-ɛ	&	{\hr}un-e	&	\cg{(rem-e)}	&	{\hr}pant-i	&	{\hr}ramn-i	&	{\hr}ann-e\\
Imroing	&	\cg{(lianti)}\su{ǁ}	&	{\hr}unn-ɛ	&	\cg{(rem-e)}	&	{\hr}lolpatl-e	&	{\hr}ram-e	&	{\hr}ann-e\\
		\lspbottomrule
	\end{tabular}
		\begin{tablenotes}\tnsize
			\item[†]	{Makuva \it{ane} ‘wind’,
								\it{huna} ‘moon’ and \it{ʤana} ‘path’ are probably from intermediate
								**anne, **hulna, and **ʤalna with simplification of the medial clusters **nn/**ln > \it{n}.
								Several languages also have assimilation of /l/ → /n/ /l{\gap}.}
			\item[‡]	{East Damer from Wulur village \it{panas} `hot' may be a loan from Malay.}
			\item[Δ]	{Emplawas has *ll > \it{s}.
								Thus, *haŋin > **aŋni > **anni > **alli > \it{usy-o},
								and *zalan > **alna > *alla > \it{as-o} ‘path, door’.}
			\item[{ǁ}] {Imroing \it{lianti} is a borrowing from Wetan.}
		\end{tablenotes}
	\end{threeparttable}
\end{adjustbox}
\end{table}

This “metathesis” is actually the result of a number of changes
which have been discussed by \citet[541ff]{BlevinsGarrett1998}
for Leti in the context of the development of productive metathesis.
The changes according to this account are summarised in \qf{04-12}
below.

{\wg\mbox{}\\\noindent\st{3.57pt}\begin{tabular}{>{\hspace{-\tabcolsep}}l<{\hspace{-\tabcolsep}}>{\hspace{-\tabcolsep}\enspace}llllllll}
\tbx{04-12}	&	\mc{8}{l}{Changes affecting consonant-final stems in \tsc{Southwest Maluku}}\\
						&a.	&*C{\#} > Cə{\#}																				&{\hr}*laŋit&>&**laŋitə	&{\hr}*panas&>&**panasə	\\
						&b.	&*ə > V\sub{[\tsc{+high}]} /V\sub{[\tsc{+high}]}C{\gap}	&**laŋitə		&>&**laŋiti	&						&	&					\\
						&c.	&*V > {\0} /ˈVC{\gap}																		&**laŋiti		&>&**laŋti	&**panasə		&>&**pansə	\\
						&d.	&*ə > a /{\gap}{\#}																			&						& &					&**pansə		&>&**pansa\\
\end{tabular}\mbox{}\\}

The first stage was epenthesis of schwa after word-final consonants.
This schwa then assimilated to the quality of a previous high
vowel resulting in *ə > \it{i} or *ə > \it{u}.
Unstressed non-final vowels were then usually lost.\footnote{
		In Luang the original final vowels are retained as short schwa [ə̆] between certain consonants.
		This schwa is perceived by speakers to be non-distinctive
		and is analysed as synchronically epenthetic \citep[6]{TaberTaber2015}.
		One example is Luang /pahna/ → [ˈpahə̆nə] ‘hot’.}
Most languages
then have final *ə > \it{a},
but schwa is retained unchanged (after non-high vowels) in Luang.
Synchronically, \citet{TaberTaber2015} analyse final schwa as
an allophone of /a/.

All languages in the \tsc{Southwest Maluku} node have evidence
that they have undergone the changes in \qf{04-12} in words of the shape VCVC{\#}.
Thus, this collection of changes can be assigned to a single
common ancestor from which the \tsc{Southwest Maluku} languages
descend and provides good evidence for the unity of this node.
However, there are several factors which complicate what is otherwise
a fairly convincing case of exclusively shared complex innovation.

Firstly, most of the \tsc{Babar} languages have lost many final vowels.
Thus, a word such as Emplawas \it{wark-o} ‘west’ could
be derived directly from intermediate *habaRat > **warat > **warak
simply with deletion of the final vowel.
However, analogy with words such as
*kulit `skin' > **ulti > \it{ulky-o} `pig skin',
*laŋit > **lanti > **lalti > \it{lalky-o} ‘sky’, and
*timuR > **timru > **kimri > \it{kipry-o} ‘east’
which show metathesis suggests that loss of final
non-high vowels in the \tsc{Babar} languages
occurred after the changes in \qf{04-12}.

Secondly, the changes in \qf{04-12} probably originally occurred only
in certain prosodic contexts.
This has resulted in some languages having productive shape alternations
partly dependent on the phonological environment in which the
words occur, though these shape alternations have also acquired purely grammatical
uses in some languages.
Thus, for instance Leti ‘black’,
from PMP *maqitəm, has three forms: \it{mɛtam} has none of the
changes in \qf{04-12}, \it{mɛtma} has all relevant changes in \qf{04-12},
and \it{mɛtm} only has the third change in \qf{04-12}.
Similarly, Leti ‘skin’ from PMP *kulit has four forms: \it{ulit,
ulti, ulty,} or \it{ult} each of which have undergone the changes
in \qf{04-12} to different extents.

Finally, for most languages of the \tsc{Southwest Maluku} node
we only have access to data from survey word lists.
A particular form for a particular language may or may not display
the changes in \qf{04-12} for reasons which cannot be determined based
on the limited data available.
Thus, for instance, PMP *laŋit > Wulur East Damer \it{lianti}
‘sky’, and *bulan → \it{wulno} ‘moon’ show that the changes in \qf{04-12}
have occurred in Wulur East Damer.
However, *haŋin > Wulur East Damer \it{anin} ‘wind’ does not show these changes.
Nonetheless, for Ilih East Damer both \it{anni,}
and \it{anin} ‘wind’ are attested.
The unmetathesised forms occurring
in the phrases \it{anin timru} ‘east wind’
and \it{anin batro} ‘west wind’,
it is thus almost certain that Wulur East Damer also has the form *[anni] ‘wind’.

To summarise, the fact that the changes in \qf{04-12} only occurred
in certain environments means that we cannot simply assign a
single change such as *VC > CV /{\gap}{\#} to Proto-Southwest Maluku.
Furthermore, the scarcity of data for many languages means that
we do not have evidence of all these changes in all available
words for every language.
However, every language of the \tsc{Southwest Maluku} node shows
evidence of the changes in \qf{04-12} having occurred in words of the shape VCVC{\#}.
We thus consider \tsc{Southwest Maluku} to be well supported.
However, more data, particularly grammatical descriptions of
more languages, would either significantly strengthen or disprove the case for \tsc{Southwest Maluku}.

Within \tsc{Southwest Maluku} we can identify five further nodes
on the basis of shared phonological innovations: the Makuva isolate
(\srf{04-06-01}), \tsc{Kisar-Luang} (\srf{04-06-02}), the East Damer isolate (\srf{04-06-03})
\tsc{Teun-Nila-Serua} (\srf{04-06-04}), and \tsc{Babar} (\srf{04-06-05}).
There are four other changes which affect several \tsc{Southwest
Maluku} languages and may provide evidence for linking some of these nodes.

Firstly, \citet{Zobel2018} posits PMP *z > **c [ʧ] for \tsc{Southwest
Maluku}, with subsequent **c > \it{s} in \tsc{Teun-Nila-Serua}
and **c > **h in \tsc{Babar}.
However, Makuva has *z = \it{ʤ}
and it is difficult to account for East Damer *z > \it{r} on
the basis of intermediate **c.
Nonetheless, we could identify a node within \tsc{Southwest Maluku}
which excludes Makuva and East Damer on the basis of shared *z > **ʧ,
with subsequent \tsc{Teun-Nila-Serua} **ʧ > \it{s}
and subsequent \tsc{Babar} **ʧ > (**s >) \it{h}.
In the absence of additional evidence we are hesitant to propose
such a node, but believe that it deserves further investigation if more data becomes available.
We also note that within the region Galolen
and \tsc{Wetar} (\srf{04-05}), as well as \tsc{Southwest Tanimbar}
have similar *z > \it{s}.

Secondly, \tsc{Southwest Maluku} languages, apart from Makuva,
usually have a partial or full merger of *d/*R/*j > \it{r}.
However, in several branches the merger is not complete and/or
there are additional complications which make it difficult to
assign a merger of any set of these segments to a higher level branch.
Thus, for instance, Teun (\srf{04-06-04})
and \tsc{Babar} (\srf{04-06-05}) apparently merge *R/*j with
a subsequent split; *R/*j > {\0},
\it{r}, but have regular *d > \it{r}.
\tsc{Kisar-Luang} regularly merges *d/*R > \it{r},
while *j has reflexes in addition to its usual *j > \it{r}.

Thirdly, all Southwest Maluku languages appear to have word-final *-aw/*a > **a.
However, data is scare for East Damer,
partly due to lack of data
and partly due to subsequent **a > \it{o} {\tl} {\0} /{\gap}{\#}.
Similarly, \tsc{Babar} normally has loss of final non-high vowels
making it difficult to discern if *-aw merged with *a,
or had some other intermediate form such as *-aw > **o.

Finally, all \tsc{Southwest Maluku} languages,
except one variety of East Damer,
have lenition of *b which may be attributable to a common ancestor,
though the diversity of reflexes: [β],
[β̞], [v], [w], [ʋ],
[f], [h] indicates that it was probably parallel development
and/or spread by diffusion.
A wave diagram representing the different sound changes affecting
Southwest Maluku languages is given in \frf{04-04} (Dotted lines
indicate changes that do not unambiguously affect all languages.)

\begin{figure}
	\includegraphics[width=\columnwidth]{SWMaluku-wave.pdf}
	\fcp{\tsc{Southwest Maluku} wave diagram}{04-04}
\end{figure}
